\documentclass{book}
\usepackage{amsmath, amssymb, amsthm}
\usepackage{geometry}
\usepackage[CJKspace]{xeCJK}  % 한글 지원
\setmainfont{Times New Roman}
\setCJKmainfont{Malgun Gothic}  % 시스템에 따라 다른 한글 글꼴 사용 가능

\geometry{
 a4paper,
 left=10mm,
 right=10mm,
 top=15mm,
 bottom=15mm
 }
 
 \setlength{\parindent}{0pt}

\title{FNN}
\author{사서림}
\date{2025.08.26.}

\begin{document}
\maketitle

\tableofcontents

\chapter{FNNv1 - Fourier Transform-based Neural Network}

내가 만든 AI모델인 FNN에 대해서 설명한다. 기본 개념 정의는 아래와 같다.

\[
S :=\text{급수 항의 개수}~~~~ C:=\text{셈 측도} ~~~~ V \in \mathbb{R}^n\text{인 입력 벡터} ~~~~ W \in \mathbb{R}^m\text{인 출력 벡터}
\]\\

이 때 벡터 V의 차원을 S로 패딩하며 이를 X라고 하자. 이때 패딩되는 값은 mean=0에 std=$10^{-7}$인 Normal분포이다. \\

다음으로 $W$의 차원별 요소 $W_j$를 구해보자. 만약 m차원일 시, 각 $W_j$는 고유한 학습 파라미터 $\{A_{sij}, A_{cij}, k_{ij}, h_{sij}, h_{cij}, b_j\}_{i=0}^{S-1}$를 가지며 다음과 같이 구할 수 있다.

\[
\omega:=\frac{2 \pi}{C \times m}
\]
\[
%W_j = \sum_{i=0}^{S-1} \left ( A_{sij} X_i \sin \left ( k_{ij} \omega i + h_{sij} X_i \right ) \right ) +\sum_{i=0}^{S-1} \left ( A_{cij} X_i \cos \left ( k_{ij} \omega i + h_{cij} X_i \right ) \right ) + b_j + \lVert V \rVert
W_j = \sum_{i=1}^{S}  \left [ A_{sij} X_i \sin \left ( \left [ \frac{k_{ij} \omega i}{\lVert V \rVert} + h_{sij} \right ] X_i \right ) \right ] +\sum_{i=1}^{S} \left [ A_{cij} X_i \cos \left ( \left [ \frac{k_{ij} \omega i}{\lVert V \rVert} + h_{cij} \right ] X_i \right ) \right ] + b_j + \lVert V \rVert
\]
\\
이 때 $A_{sij}, A_{cij}, k_{ij}, h_{sij}, h_{cij}, b_j$는 학습해야 하는 파라미터이다. 이때 초기화는 $A, k, h$ 계열은 SIREN초기화로 $\text{fan\_in}:=S$로 하여 초기화 하고 $b_j$는 mean=0, std=0.001로 초기화 한다.\\

최종적으로 FNN은 다음과 같이 표현할 수 있다. (MLP Header를 거치지 않는 단층 FNN을 SLF라고 부르기도 한다.)
\[
\text{FNN}:V \rightarrow W
\]

그리고 성능 향상을 위하여 FNN의 결과인 fnn에 대하여 MLP을 거친다. 
\[
L_1:\mathbb{R}^m\rightarrow \mathbb{R}^{8m}~~~~L_2:\mathbb{R}^{8m}\rightarrow \mathbb{R}^m
\]
\[
M(fnn) = L_2 \circ \text{GELU} \circ L_1(\text{fnn}) 
\]

\chapter{Hyper-FNN - Hypernetwork-based FNN}

기존 FNN 모델은 입력 $V$에 관계없이 고정된 학습 파라미터 $\{A_{sij}, A_{cij}, k_{ij}, h_{sij}, h_{cij}\}$를 사용한다. 이는 모델의 표현력을 제한할 수 있다. 이러한 한계를 극복하기 위하여, FNN의 핵심 파라미터들을 입력 벡터 $V$에 따라 동적으로 생성하는 하이퍼네트워크(Hypernetwork) 구조를 도입한 Hyper-FNN을 제안한다.

\section{하이퍼네트워크의 정의}

하이퍼네트워크 $H$는 FNN의 학습 파라미터 전체를 생성하는 작은 신경망이다. $H$는 입력 벡터 $V \in \mathbb{R}^n$을 입력으로 받아, FNN의 모든 파라미터 $\{A, k, h\}$를 포함하는 거대한 파라미터 텐서 $P \in \mathbb{R}^{m \times 5 \times S}$를 출력한다.

\[
H: V \in \mathbb{R}^n \rightarrow P(V) \in \mathbb{R}^{m \times 5 \times S}
\]
\\
이때 출력 텐서 $P(V)$는 다음과 같이 개별 파라미터 함수들로 분해하여 생각할 수 있다.
\[
P(V) := \{A_{sij}(V), A_{cij}(V), k_{ij}(V), h_{sij}(V), h_{cij}(V)\}_{j=0, i=0}^{m-1, S-1}
\]
\\
여기서 각 파라미터 $A_{sij}(V), k_{ij}(V)$ 등은 더 이상 상수가 아니라, 입력 $V$에 대한 함수가 된다. 하이퍼네트워크 $H$ 자체는 일반적으로 간단한 다층 퍼셉트론(MLP)으로 구성되며, $H$의 학습 파라미터들이 실제 Hyper-FNN에서 학습해야 할 대상이 된다.

\section{Hyper-FNN의 순전파 수식}

하이퍼네트워크 $H$를 통해 생성된 동적 파라미터들을 이용하여, 출력 벡터 $W$의 각 요소 $W_j$는 다음과 같이 새롭게 정의된다. 기존 FNN 수식의 모든 파라미터가 $V$에 대한 함수로 대체된 것을 볼 수 있다.

\[
\omega:=\frac{2 \pi}{C \times m}
\]
\[
W_j(V) = \sum_{i=1}^{S} \left [ A_{sij}(V) X_i \sin \left ( \left [ \frac{k_{ij}(V) \omega i}{\lVert V \rVert} + h_{sij}(V) \right ] X_i \right ) \right ]
\]
\[
+\sum_{i=1}^{S} \left [ A_{cij}(V) X_i \cos \left ( \left [ \frac{k_{ij}(V) \omega i}{\lVert V \rVert} + h_{cij}(V) \right ] X_i \right ) \right ] + b_j + \lVert V \rVert
\]
\\
여기서 $b_j$는 기존과 같이 고정된 학습 파라미터로 남겨둘 수도 있고, 하이퍼네트워크가 함께 생성하도록 설계할 수도 있다.

\section{Hyper-FNN의 의의}

Hyper-FNN은 입력 데이터에 따라 완전히 맞춤화된 푸리에 급수 근사를 수행하는 `메타 모델(Meta-model)'로 볼 수 있다. 이 구조는 다음과 같은 장점을 가진다.

\begin{itemize}
    \item \textbf{입력 적응성(Input Adaptability):} 모든 입력에 대해 최적화된 파라미터를 동적으로 생성하므로, 훨씬 더 복잡하고 다양한 데이터 분포를 모델링할 수 있는 잠재력을 가진다.
    \item \textbf{파라미터 효율성(Parameter Efficiency):} 거대한 파라미터 텐서(크기: $m \times 5 \times S$)를 직접 저장하는 대신, 그보다 훨씬 작은 하이퍼네트워크 $H$의 파라미터만 저장하면 되므로, 전체 모델의 크기를 효율적으로 관리할 수 있다.
\end{itemize}

\chapter{FNN v2 - Fourier Feature Neural Network (Matrix Form)}

\section{기본 정의}
입력 $V \in \mathbb{R}^n$, 출력 $W \in \mathbb{R}^m$.
푸리에 급수 항의 개수를 $S$라 하자.

\[
\omega := \frac{2\pi}{C \cdot m}, \qquad
X = \text{Pad}(V; S)
\]

여기서 $\text{Pad}$는 길이를 $S$로 맞추기 위한 패딩 연산이다.

\section{Fourier Feature 벡터}
입력 $X \in \mathbb{R}^S$에 대하여 다음의 푸리에 특징 벡터를 정의한다.
\[
\Phi(X) =
\begin{bmatrix}
\sin(\Omega X + H_s) \\
\cos(\Omega X + H_c)
\end{bmatrix}
\in \mathbb{R}^{2S}
\]
여기서 $\Omega, H_s, H_c \in \mathbb{R}^{S \times S}$는 학습 파라미터 행렬이다.

\section{출력 수식}
FNN v2의 출력은 다음과 같이 정의된다.
\[
W = A \, \Phi(X) + b + \lVert V \rVert
\]
여기서 $A \in \mathbb{R}^{m \times 2S}$, $b \in \mathbb{R}^m$은 학습 파라미터이다.

\section{헤더 MLP}
성능 향상을 위해, $W$에 대해 2단계 MLP를 적용한다.
\[
M(W) = L_2 \circ \sigma \circ L_1(W),
\]
\[
L_1:\mathbb{R}^m \to \mathbb{R}^{8m}, \quad
L_2:\mathbb{R}^{8m} \to \mathbb{R}^m, \quad
\sigma=\text{GELU}.
\]

\section{Hyper-FNN v2 - Hypernetwork-based FNN (Matrix Form)}

\subsection{하이퍼네트워크의 정의}
하이퍼네트워크 $H$는 입력 $V \in \mathbb{R}^n$을 받아
FNN v2의 모든 파라미터 행렬을 동적으로 생성한다.
\[
H: V \mapsto \{A(V), \Omega(V), H_s(V), H_c(V), b(V)\}.
\]

\subsection{순전파}
생성된 파라미터를 이용해 출력은 다음과 같이 정의된다.
\[
W(V) = A(V) \, \Phi\!\big(X; \Omega(V), H_s(V), H_c(V)\big) + b(V) + \lVert V \rVert.
\]

\subsection{의의}
\begin{itemize}
\item \textbf{입력 적응성}: 파라미터가 $V$에 의존하므로 데이터 분포 변화에 강함.
\item \textbf{파라미터 효율성}: 거대한 파라미터를 직접 저장하지 않고, 작은 하이퍼네트워크 $H$만 학습하면 됨.
\end{itemize}


\chapter{GFN - Global Fourier Network}

이제는 한번 FNNv1의 푸리에 급수를 2차원 푸리에 급수로 확장한 식을 적어보자. 일단 그 이전에 FNNv1의 식을 보고 2차원 푸리에 급수를 보자\\

[FNNv1의 수식]
\[
W_j = \sum_{i=1}^{S}  \left [ A_{sij} X_i \sin \left ( \left [ \frac{k_{ij} \omega i}{\lVert V \rVert} + h_{sij} \right ] X_i \right ) \right ] +\sum_{i=1}^{S} \left [ A_{cij} X_i \cos \left ( \left [ \frac{k_{ij} \omega i}{\lVert V \rVert} + h_{cij} \right ] X_i \right ) \right ] + b_j + \lVert V \rVert
\]\\

[2차원 푸리에 급수 수식]
\[
C_{m,n} =\frac{1}{WH} \int_0^W \int_0^H f(x,y) e^{-2\pi i \left ( \frac{x}{W} + \frac{y}{H} \right )}
\]
\[
F(x,y)= C_{x,y} \sum_{x=0}^W \sum_{y=0}^H e^{2 \pi i \left ( \frac{x}{W} + \frac{y}{H} \right )}
\]\\

이제는 2차원 푸리에 급수와 FNNv1의 수식을 이용하여 GFN의 초안 수식을 적어보겠다.

입력 $M:=R^{n \times m}$, 출력 $N:=R^{k \times h}$인 행렬이라고 생각해보고 그리고 학습 파라미터 $z_{0}, z_{1}, z_{2} \cdots z_{n}$인 스칼라를 생각해보자. 그럼 $M_{x,y}$는 다음과 같다. 참고로 노름은 프로베니우스 노름이다

\[
N_{x,y} = \sum_{x=0}^{W-1} \sum_{y=0}^{H-1} z_{xy0N_{x,y}} M_{x,y} \Re \left ( \exp \left (2 \pi M_{x,y} i \left [ \dfrac{ \dfrac{xz_{xy1N_{x,y}}}{W} + \dfrac{yz_{xy2N_{x,y}}}{H} }{\lVert M \lVert} + z_{xy3N_{x,y}} \right ] \right ) \right ) + \lVert M \rVert
\]\\

\section{GFN v2 - Global Fourier Network (Optimized Form)}

FNNv2 아키텍처의 핵심 철학인 '특징 추출'과 '선형 결합'의 분리 개념을 2차원 이미지에 적용하여, GFNv1의 계산적 복잡성 문제를 해결한 GFN v2를 제안한다.

\subsection{기본 정의}
입력 이미지를 $M \in \mathbb{R}^{C \times H \times W}$ (C는 채널 수), 최종 출력을 $N \in \mathbb{R}^k$ (k차원 벡터, 분류 문제의 경우)라 하자.
추출할 푸리에 특징의 개수를 $K$라고 정의한다.

\subsection{전역 푸리에 특징 (Global Fourier Feature) 벡터}
GFN v2의 핵심은 이미지 전체로부터 고정된 크기의 특징 벡터 $\Psi(M) \in \mathbb{R}^{2K}$를 추출하는 단일 레이어에 있다.

$k \in \{1, 2, \dots, K\}$인 각 푸리에 특징에 대하여, 사인 성분 $s_k$와 코사인 성분 $c_k$를 FNNv1의 철학을 계승하여 정의한다:

% --- 여기가 수정된 수식입니다 ---
\[
s_k(M) = \frac{1}{C \cdot H \cdot W} \sum_{c=0}^{C-1} \sum_{x=0}^{H-1} \sum_{y=0}^{W-1} M_{c,x,y} \sin\left( \frac{u_k x + v_k y}{\lVert M \rVert_F} + h_{sk} \right)
\]
\[
c_k(M) = \frac{1}{C \cdot H \cdot W} \sum_{c=0}^{C-1} \sum_{x=0}^{H-1} \sum_{y=0}^{W-1} M_{c,x,y} \cos\left( \frac{u_k x + v_k y}{\lVert M \rVert_F} + h_{ck} \right)
\]

여기서 $\{u_k, v_k, h_{sk}, h_{ck}\}_{k=1}^{K}$는 학습 가능한 파라미터이며, $\lVert M \rVert_F$는 입력 행렬 $M$의 프로베니우스 노름(Frobenius Norm)이다.

\subsection{학습 안정성을 위한 정규화 (Normalization for Training Stability)}
위 수식에서 정규화 항 $\frac{1}{C \cdot H \cdot W}$은 매우 중요한 역할을 한다.
단순히 모든 픽셀 값을 합산(\texttt{summation})할 경우, 입력 이미지의 해상도가 커짐에 따라 특징 벡터의 크기가 비정상적으로 커져 손실 함수(Loss Function)의 폭주(explosion)를 유발할 수 있다.
따라서 전체 픽셀 수로 나누어주는 정규화 과정을 통해, 이 연산은 '총합'이 아닌 '평균'을 계산하게 된다. 
이는 특징 벡터의 스케일을 입력 이미지의 크기와 무관하게 일정 범위로 유지시켜, 안정적인 학습을 가능하게 하는 핵심적인 장치이다.

최종적으로 전역 푸리에 특징 벡터 $\Psi(M)$는 이 모든 성분을 하나로 결합한 것이다.
\[
\Psi(M) = 
\begin{bmatrix}
s_1(M) \\ c_1(M) \\ s_2(M) \\ c_2(M) \\ \vdots \\ s_K(M) \\ c_K(M)
\end{bmatrix}
\in \mathbb{R}^{2K}
\]

\subsection{출력 수식 (분류 모델의 경우)}
추출된 특징 벡터 $\Psi(M)$를 사용하여, 최종 출력 $N \in \mathbb{R}^k$는 간단한 선형 변환 및 잔차 연결로 정의된다.
\[
N = A \, \Psi(M) + b
\]
여기서 $A \in \mathbb{R}^{k \times 2K}$와 $b \in \mathbb{R}^k$는 학습 가능한 파라미터이다. 필요에 따라 이 위에 MLP 헤더를 추가할 수 있다.

(주의: FNNv1의 잔차 연결 항 $\lVert M \rVert_F$는, CrossEntropyLoss를 사용하는 분류 문제에서는 출력 로짓(logit)의 스케일을 왜곡시킬 수 있으므로 여기서는 제외한다.)

\subsection{의의 및 장점}
이 GFN v2 구조는 GFN v1의 철학을 계승하면서도, 다음과 같은 결정적인 장점을 가진다.
\begin{itemize}
    \item \textbf{계산적 실현 가능성}: 파라미터 수와 계산량이 더 이상 이미지 크기의 제곱($H^2 \times W^2$)이 아닌, 푸리에 특징의 개수 $K$에만 선형적으로 비례한다. 이로써 실제 하드웨어에서 훈련이 가능해진다.
    \item \textbf{입력 적응성}: FNNv1의 핵심 아이디어였던 노름 정규화를 통해, 각 이미지의 특성에 맞춰 동적으로 다른 주파수 분석을 수행하는 능력을 유지한다.
\end{itemize}

\chapter{GFN v2.1 - Fixed Fourier Feature Network}

이전 GFNv2 모델에 대한 실험을 통해, 공간 주파수 파라미터($u_k, v_k$)를 학습 가능하게 둘 경우, 모델이 손실(Loss)을 줄이는 가장 쉬운 경로, 즉 소수의 지배적인 저주파 특징(low-frequency features)만을 학습하는 경향이 있음을 발견했다. 이로 인해 푸리에 특징의 개수 $K$를 늘려도 모델의 표현력이 실질적으로 향상되지 않는 문제가 발생한다.

이 문제를 해결하기 위해, 각 푸리에 특징 분석기가 서로 다른 전문 분야를 갖도록 강제하는, GFN v2.1 (Fixed Fourier Feature Network) 아키텍처를 제안한다.

\section{핵심 아이디어: 고정된 주파수 기저 (Fixed Frequency Basis)}
GFN v2.1의 핵심 철학은 주파수 파라미터 $u_k, v_k$를 학습의 대상에서 제외하는 것이다. 대신, 이 파라미터들을 저주파부터 고주파까지 넓은 스펙트럼을 포괄하도록 미리 정의된 **고정된 상수**로 설정한다.

이를 통해 $K$개의 각 푸리에 특징 추출기는 각각 고유한 주파수 대역을 담당하는 '전문가'가 되어, 모델이 이미지의 다양한 측면(전체적인 윤곽부터 미세한 질감까지)을 모두 분석하도록 강제한다.

\section{전역 푸리에 특징 벡터 (수정된 정의)}
$k \in \{1, 2, \dots, K\}$인 각 푸리에 특징에 대하여, 사인 성분 $s_k$와 코사인 성분 $c_k$를 다음과 같이 재정의한다:

\[
s_k(M) = \frac{1}{C \cdot H \cdot W} \sum_{c=0}^{C-1} \sum_{x=0}^{H-1} \sum_{y=0}^{W-1} M_{c,x,y} \sin\left( \frac{\hat{u}_k x + \hat{v}_k y}{\lVert M \rVert_F} + h_{sk} \right)
\]
\[
c_k(M) = \frac{1}{C \cdot H \cdot W} \sum_{c=0}^{C-1} \sum_{x=0}^{H-1} \sum_{y=0}^{W-1} M_{c,x,y} \cos\left( \frac{\hat{u}_k x + \hat{v}_k y}{\lVert M \rVert_F} + h_{ck} \right)
\]

여기서 가장 중요한 변화는 $\hat{u}_k, \hat{v}_k$가 더 이상 학습 가능한 파라미터(\texttt{nn.Parameter})가 아닌, 미리 정의된 상수(fixed constant)라는 점이다. 이제 이 레이어에서 학습되는 파라미터는 오직 위상 변이를 담당하는 $\{h_{sk}, h_{ck}\}_{k=1}^{K}$ 뿐이다.

$\hat{u}_k, \hat{v}_k$는 가우시안 분포 $\mathcal{N}(0, \sigma^2)$에서 샘플링하거나, 로그 선형 공간(log-linear space)에 걸쳐 균등하게 설정하는 등, 다양한 방식으로 미리 정의될 수 있다.

\section{의의 및 장점}
이 '고정된 주파수' 접근법은 GFN 아키텍처에 다음과 같은 결정적인 이점을 제공한다.

\begin{itemize}
    \item \textbf{특징의 다양성 보장}: 모델이 '게으른' 최적화 경로를 선택하는 것을 원천적으로 차단하고, 저주파(전체 구조)와 고주파(디테일, 질감) 특징을 모두 분석하도록 강제하여 모델의 표현력을 극대화한다.
    \item \textbf{학습 안정성 및 속도 향상}: 학습해야 할 파라미터의 수를 줄이고, 손실 공간(Loss Landscape)을 더 단순하게 만들어 학습을 더 빠르고 안정적으로 만든다.
    \item \textbf{진정한 계층적 학습}: GFN 레이어가 다양한 스케일의 특징을 안정적으로 추출해주기 때문에, 그 뒤에 오는 MLP 헤더는 이 풍부한 정보들을 조합하여 더 의미 있는 고차원적인 추론을 하는 데 집중할 수 있게 된다.
\end{itemize}

이것은 "Fourier Features Let Networks Learn High Frequency Functions"와 같은 선행 연구에서 영감을 얻은, 매우 강력하고 검증된 접근 방식이다.


\end{document}